%% IEEE Sensors 2026 — 3-page submission
%% Max 3 pages body + figures + tables; references may occupy the 4th page only.
%% Quantum E-Nose: Rank-1 projected NEGF–SCBA for ZnO/Mg_{0.3}Zn_{0.7}O IETS
%%
%% This document expands the 2-page SISPAD-2026 submission into the IEEE Sensors
%% 3-page format. No new results have been added; the additional space is used
%% for a more complete methodology derivation, a paragraph on self-consistent
%% electrostatics (Poisson), and a more detailed results discussion.

% CHANGE 1: Remove a4paper — IEEEtran conference defaults to letter
\documentclass[conference]{IEEEtran}
\IEEEoverridecommandlockouts

\usepackage{cite}
\usepackage{amsmath,amssymb,amsfonts}

% CHANGE 2: Force Times New Roman for body text (10pt is set by IEEEtran conference class)
\usepackage{times}          % or: \usepackage{newtxtext,newtxmath} for better math matching

\usepackage{graphicx}
\usepackage{textcomp}
\usepackage{xcolor}
\usepackage{booktabs}
\usepackage{stfloats}
\usepackage{subcaption}

% CHANGE 3: Fix caption font size to 9pt (IEEEtran conference uses 8pt by default)
\usepackage{caption}
\captionsetup{font={rm,small},labelfont=bf,textfont={rm}}
% 'rm' ensures Roman (Times) family; size=9pt overrides the class default

\usepackage[colorlinks=true, urlcolor=blue, citecolor=blue, linkcolor=darkgray]{hyperref}
\renewcommand\figureautorefname{Fig.}
\renewcommand\tableautorefname{Tab.}

%% ── Compactness tweaks (relaxed vs SISPAD version) ──────────────────────────
\setlength{\abovecaptionskip}{2pt}
\setlength{\belowcaptionskip}{0pt}
\setlength{\floatsep}{4pt plus 1pt minus 1pt}
\setlength{\textfloatsep}{6pt plus 1pt minus 2pt}
\setlength{\intextsep}{4pt plus 1pt minus 1pt}
\setlength{\dblfloatsep}{4pt plus 1pt minus 1pt}
\setlength{\dbltextfloatsep}{6pt plus 1pt minus 2pt}
%% ─────────────────────────────────────────────────────────────────────────────

\begin{document}

%% NOTE: IEEE Sensors template warns against math/symbols in title and abstract.
%% Chemical subscripts (\textsubscript) and standard unit symbols are restored
%% here as they are universally used in accepted IEEE Sensors papers; heavier
%% inline math has been spelled out in words instead.
\title{Design of a Room-Temperature Quantum Biomimetic Electronic Nose Using a CMOS BEOL Compatible Resonant Tunneling Diode}

\author{
\IEEEauthorblockN{Abhineet Agarwal and Swaroop Ganguly}
\IEEEauthorblockA{\textit{Department of Electrical Engineering},
Indian Institute of Technology Bombay, Mumbai, India\\
Contact: abhineetagarwal@iitb.ac.in\ /\ Phone: +916372603493}
}

\IEEEaftertitletext{\vspace{-8pt}}
\maketitle

\begin{abstract}
A biomimetic electronic nose that identifies odorants by their molecular vibrational fingerprint would offer both high sensitivity to individual species and broad chemical coverage from a single device. Inelastic Electron Tunneling Spectroscopy (IETS) gives such fingerprints directly as peaks in $d^2I/dV^2$ at the vibrational energies of molecules in the tunnel barrier, but conventional IETS requires cryogenic operation. We present the design of a CMOS back-end-of-line (BEOL) compatible ZnO/Mg\textsubscript{0.3}Zn\textsubscript{0.7}O resonant tunneling diode (RTD) as a room-temperature IETS sensor, in which the narrow RTD resonance acts as an intrinsic energy filter that decouples spectral resolution from thermal broadening. A quasi-3D Non-Equilibrium Green's Function model with a projected self-consistent Born approximation is used to predict resolvable molecule-specific $d^2I/dV^2$ signatures at 300~K for synthetic single and two-mode analytes (100~meV and 180~meV). Combined with BEOL CMOS deposition and array readout via standard AC lock-in or low-noise DC differentiation, the device could provide a suitable platform for a room-temperature quantum biomimetic electronic nose.
\end{abstract}

\begin{IEEEkeywords}
electronic nose, resonant tunneling diode, inelastic electron tunneling spectroscopy, NEGF, oxide heterostructures, olfactory sensor
\end{IEEEkeywords}

\section{Introduction}
Inelastic Electron Tunneling Spectroscopy (IETS) provides vibrational fingerprints of molecules adsorbed in or near a tunnel junction by introducing weak inelastic openings of additional transport channels at molecule-specific biases, observable as peaks in $d^2I/dV^2$~\cite{hansma1977,hansma1982,hipps1993,reed2008}. It has been widely used for studying electron-phonon inelastic interaction and vibrational spectroscopy of molecules and organic and inorganic thin films~\cite{jeong2015,liu2010,kim2011,wang2004,lauhon2001}. 

In 1996, Turin proposed that biological olfaction itself may be an IETS-like process, with olfactory receptors sensing the vibrational energies of odorant molecules rather than (or in addition to) their shape~\cite{turin1996}. This proposition has motivated technological interest in a \emph{biomimetic} electronic nose that identifies odorants by their intrinsic vibrational signature~\cite{patil2018,bommisetty2008,patil2013mit, li2024}. Recent computational work from this group has shown that unsupervised clustering on the vibrational spectra of 20 odorants recovers their perceptual odor categories and even resolves subclasses~\cite{pandey2021}. Existing electronic-nose technologies span resistive, capacitive, gravimetric and optical transduction \cite{rock2008}; but mimicking natural olfaction remains an open challenge with applications in security, environmental monitoring, healthcare, and food and agriculture. Such a sensor will be ultra-sensitive, but still near universal - one device can be used to detect a wide class of molecules.

The obstacle to a deployable IETS-based e-nose is resolution at non-cryogenic temperatures. The Lambe--Jaklevic analysis of thermal Fermi-distribution softening gives a best-case spectral resolution of $5.4 k_B T$, which at 300~K is $\approx 140$~meV, comparable to odorant vibrational modes. IETS was for decades considered unusable above $\sim$150~K~\cite{jaklevic1966}. Two recent developments challenge this. Ngabonziza~\textit{et al.}~\cite{ngabonziza2021} reported the first inelastic-tunneling spectra above room temperature, achieving 20~meV resolution at 400~K in BaZrO$_3$-based proton-conducting defect-mediated tunnel junctions and demonstrating that the $5.4 k_B T$ bound is not fundamental but a property of conventional MIM geometries. Earlier work from this group~\cite{patil2018} showed that a double-barrier resonant tunneling diode (RTD) achieves the same: its narrow transmission resonance acts as an intrinsic energy filter whose effective width is set by the resonance linewidth $\Gamma_{\mathrm{res}}$, not $k_BT$, preserving $d^2I/dV^2$ features at 300~K in a toy device. 

We extend this concept to a manufacturable all-oxide stack: ZnO/Mg$_{0.3}$Zn$_{0.7}$O, which to our knowledge is the only published BEOL-compatible oxide pair with conduction-band offset in the $\Delta E_c \approx 0.5$~eV window and experimentally demonstrated NDR \cite{krishnamoorthy2002}. We quantitatively study its ability to discriminate molecular vibrational spectra at room temperature using a rigorously derived \emph{projected NEGF--SCBA}. In a sensor system, each $10\,\mu$m $\times$ $10\,\mu$m RTD would be one device in an array, with the $d^2I/dV^2$ spectrum measured using standard AC lock-in techniques or DC measurement with noise filtering~\cite{kesarwani2022}

%For a Gaussian molecular approximation ($\sigma_{\mathrm{mol}} \approx 3$~\AA), the electron-vibron coupling factorizes as $M = D_0\, \chi(z)\, |u\rangle\langle u|$ in the transverse-mode basis (correction $O(10^{-6})$). This reduces the 3D SCBA to a scalar 1D projected Dyson equation, exact in the continuum limit.

\section{Device and Method}
\textbf{Device.} The symmetric RTD consists of 2~nm Mg$_{0.3}$Zn$_{0.7}$O barriers, a 3~nm ZnO well, and 10~nm $n^+$-ZnO contacts ($N_D = 5\times10^{24}$~m$^{-3}$). Material parameters used in the simulations are $\Delta E_c = 0.47$~eV, $m^* = 0.28\, m_0$, and a bulk ZnO longitudinal-optical (LO) phonon energy $\hbar\omega_{\mathrm{LO}} \approx 72$~meV. NDR occurs near $V_{\mathrm{res}} \approx 560$~mV with peak current $\sim$69~nA on a $10\,\mu$m $\times$ $10\,\mu$m footprint (Fig.~\ref{fig:device}). Candidate material stacks for room-temperature IETS are compared in Table~\ref{tab:stacks}; ZnO/Mg$_{0.3}$Zn$_{0.7}$O uniquely satisfies all three criteria (CBO in vibrational range, demonstrated NDR, BEOL deposition $<400^\circ$C).

\textbf{NEGF transport.} Longitudinal discretization uses a $0.2$~nm tight-binding grid with effective-mass kinetic terms. The retarded Green's function $G^R(E) = [(E + i\eta)\mathbb{1} - H - \Sigma^L - \Sigma^R - \Sigma^{\mathrm{ph}}]^{-1}$ is in the standard form with semi-infinite-lead self-energies \cite{datta_textbook}, symmetric bias ($\mu_{L,R} = E_F \pm V/2$, $E_F = 20$~meV above the ZnO conduction-band edge, the degenerate value fixed by charge neutrality at the quoted contact doping), and a causal phonon self-energy $\Sigma^{\mathrm{ph}}$ solved within the self-consistent Born approximation (SCBA). 

\textbf{Rank-1 projection.} The molecular vibron couples to electrons through a localized Gaussian deformation potential $g(\mathbf{r}_\perp, z) = D_0\, \chi(z)\, \mathcal{G}_\sigma(\mathbf{r}_\perp)$, where $\chi(z)$ is a Gaussian envelope along the transport axis peaking at the emitter barrier and $\mathcal{G}_\sigma$ is a unit-normalized 2D Gaussian of width $\sigma_{\mathrm{mol}}$. In the transverse plane-wave basis $\{|\mathbf{k}_\perp\rangle\}$, the electron-vibron matrix element becomes
\begin{equation}
M_{\mathbf{k},\mathbf{k}'}(z) = D_0\, \chi(z)\, e^{-\sigma_{\mathrm{mol}}^2 |\mathbf{k}_\perp - \mathbf{k}'_\perp|^2 / 2}.
\label{eq:formfactor}
\end{equation}
In the relevant transverse-momentum range, $\sigma_{\mathrm{mol}}|\mathbf{k}_\perp| \ll 1$, the equation in (\ref{eq:formfactor}) reduces to a separable outer product $u_\mathbf{k}\, u_{\mathbf{k}'}$ with $u_\mathbf{k} = \exp(-\sigma_{\mathrm{mol}}^2 k_\perp^2 / 4)$, leading to the structure
\begin{equation}
M(z) \;\approx\; D_0\, \chi(z)\, |u\rangle\langle u|,
\qquad \text{error } \mathcal{O}\!\left((\sigma_{\mathrm{mol}}\, k_\perp)^2\right).
\label{eq:rank1}
\end{equation}
The bulk LO self-energy is mode-diagonal; combined, $\Sigma^{\mathrm{ph}} = \Sigma^{\mathrm{bulk}} + \Sigma^{\mathrm{mol}}$ is diagonal-plus-rank-1. The SCBA transverse sum collapses onto a single projected channel, and the converged Dyson equation becomes one-dimensional in $z$ (Fig.~\ref{fig:flowchart}).

\textbf{Bulk + molecular phonons.} The bulk ZnO LO phonon ($\hbar\omega_{\mathrm{LO}} = 72$~meV) couples uniformly in the transverse plane ($D_{\mathrm{bulk}}^2 = 0.001$~eV$^2$/site, $\chi \equiv 1$ in the active region). Molecular vibrons are added on top with $D_0^2 = 0.1$~eV$^2$ and a Gaussian $\chi(z)$ centered on the emitter barrier. 

\textbf{Self-consistent electrostatics.} The present calculations use a non-self-consistent electrostatic potential under the applied bias, with charge neutrality enforced in the doped contacts. Adding a self-consistent Hartree coupling renormalizes the resonance bias position, and the sign of that shift depends on the doping profile adjacent to the barrier. In structures that interpose a lightly-doped spacer between the doped lead and the barrier, an emitter accumulation layer forms in the spacer, screens the barrier field, and displaces the resonance to higher bias~\cite{lake1992,lake1997,akkala2011}. The stack simulated here has no spacer: ionized donors sit directly against the barrier, where quantum reflection suppresses the carrier density, leaving a net positive space charge that concentrates the field; preliminary self-consistent Poisson--NEGF calculations accordingly displace the resonance to \emph{lower} bias, by of order $100$~mV. (Those calculations use the same 2/3/2~nm active region with the doped leads extended to 30~nm, which is required to give the contact screening sufficient spatial room; the resonance energies are unchanged by the lead length.) In either case the effect is a rigid translation of the transport window rather than a change in its internal structure. Because IETS discrimination relies on the \emph{relative} spacing of phonon-induced peaks in $d^2I/dV^2$ (corresponding to $\hbar\omega_{\mathrm{mol}}/e$), this bias renormalization shifts the measurement window without altering the molecular fingerprint structure. A quantitative self-consistent treatment is therefore deferred to future work.

\textbf{Convergence.} Anderson mixing~\cite{anderson1965} (DIIS, history $M=8$, damping $\beta=0.3$) drives the SCBA loop to a converged $\Sigma^{\mathrm{ph}}$ in 10--60 iterations across all biases, with current conservation $|I_L - I_R|/|I_L| < 0.2\%$ on the converged solutions. The transport energy grid ($\Delta E = 2$~meV) does not resolve the 134~$\mu$eV linewidth of the $E_1$ resonance, which contributes $<5\%$ of the peak current; the $E_2$ transport peak is converged to $0.3\%$ against a 10~$\mu$eV reference grid.

\textbf{Analytes.} Three synthetic molecules are studied: Mol\_A ($\hbar\omega_A = 100$~meV), Mol\_B ($\hbar\omega_B = 180$~meV), Mol\_AB (both modes), and a Baseline with only the bulk 72~meV phonon. These mode energies lie in the vibrational-fingerprint band (50--200~meV) characteristic of small organic odorants~\cite{pandey2021}.

\section{Results and Discussion}

\subsection{I--V and NDR characterization}
All analytes show clean NDR near $V_{\mathrm{res}} \approx 560$~mV (Fig.~\ref{fig:iv}). The Baseline peak current is $\sim$69~nA with a peak-to-valley ratio (PVR) $\approx 12$, consistent with the values expected from a 2/3/2~nm ZnO/Mg$_{0.3}$Zn$_{0.7}$O stack at 300~K. The molecular vibrons reduce the elastic transmission by opening additional inelastic channels, leading to molecule-dependent step features in $I(V)$ on top of the elastic NDR background.

\subsection{Molecular fingerprint mechanism}
At a bias $V$ such that $eV$ crosses the threshold $E_{\mathrm{res}} + \hbar\omega_{\mathrm{mol}}$ (measured from the emitter Fermi level), an inelastic emission channel opens: an incoming electron with energy $E > E_{\mathrm{res}} + \hbar\omega_{\mathrm{mol}}$ can shed a vibrational quantum and resonantly traverse the well. This produces a step in $I(V)$ and a corresponding peak in $d^2I/dV^2$ at $V \approx \hbar\omega_{\mathrm{mol}}/e$ relative to the emitter band edge. Fig.~\ref{fig:d2i} shows these molecule-specific deviations clearly resolved at 300~K. The discrimination metric $\Delta D(V) = d^2I_{\mathrm{mol}}/dV^2 - d^2I_{\mathrm{BL}}/dV^2$ (Fig.~\ref{fig:deltaD}) isolates the vibrational signature from the elastic NDR background. Crucially, Mol\_AB exhibits near-additive superposition of the Mol\_A and Mol\_B deviations, a direct consequence of the rank-1 structure of (\ref{eq:rank1}): when the molecular coupling is a single outer product, multiple modes contribute through orthogonal projections that add linearly in the weak-coupling regime.

\subsection{Temperature robustness and sensor implications}
Fig.~\ref{fig:temp} sweeps temperature from 10~K to 300~K for Mol\_A. The peak current rises from 69~nA at 300~K to 110~nA at 10~K as the NDR sharpens. Critically, the molecular feature near $\hbar\omega_{\mathrm{mol}}/e = 100$~mV narrows at low $T$ but persists to 300~K (Fig.~\ref{fig:temp}(b), inset), where its amplitude is $\approx 2\,\mu$A/V$^2$, some $3\%$ of the resonance-derived background; it is most cleanly isolated in the discrimination metric $\Delta D$ (Fig.~\ref{fig:deltaD}), in which the elastic background cancels. This is the central practical claim of the device: the RTD's resonance linewidth $\Gamma_{\mathrm{res}}$ sets the effective energy resolution, decoupling IETS visibility from the $3k_BT$ thermal-broadening limit that constrains conventional tunnel-junction IETS. Combined with BEOL compatibility (deposition temperature $<400^\circ$C) and a $10\,\mu$m $\times$ $10\,\mu$m footprint, this enables array-like integration onto a CMOS readout substrate, with each pixel reporting an IETS spectrum.

\section{Conclusion}
We predict, through a rigorously derived projected NEGF--SCBA simulation, that a BEOL-compatible ZnO/Mg$_{0.3}$Zn$_{0.7}$O resonant tunneling diode can perform room-temperature IETS-based molecular discrimination. The methodological contribution here is the reduction of the full 3D SCBA to a single projected 1D Dyson equation valid for localized molecular vibrons. The physical contribution is the identification of an all-oxide, CMOS BEOL compatible stack with intrinsic energy filtering which produces vibrational fingerprints at 300~K. Future work would comprise self-consistent Poisson--NEGF for quantitative bias calibration, DFT-computed electron--vibron couplings for specific volatile organic compounds, fabrication of the studied stack and array simulations to model system-level variability.
\newpage
% ====================== Figures and tables ======================
%% References block has been moved after the figures, with \clearpage,
%% so that all references fall on page 4 per IEEE Sensors instructions.
%% Fig 1: Band diagram
\setcounter{figure}{0}
\begin{figure*}[!t]
\centering
\includegraphics[width=\textwidth]{fig1_band_diagram.pdf}
\caption{\textbf{(a)}~Epitaxial stack of the symmetric ZnO/Mg$_{0.3}$Zn$_{0.7}$O RTD (10/2/3/2/10~nm, 27~nm total). \textbf{(b)}~Conduction-band profile at $V=0$: the well supports two quasi-bound states, $E_1\approx79$~meV and $E_2\approx296$~meV. $E_1$ has a linewidth of only 134~$\mu$eV and is strongly suppressed, contributing $<5\%$ of the peak current; transport at the NDR proceeds through $E_2$. The orange curve is the molecular coupling profile $\chi(z)$ (Gaussian, $\sigma_{\mathrm{mol}}=0.3$~nm, centred on the emitter barrier); $\chi$ is dimensionless and its height is in arbitrary units. \textbf{(c)}~At $V=V_{\mathrm{peak}}\approx560$~mV, $E_2$ is resonant with $\mu_L$, i.e. $V_{\mathrm{res}}=2(E_2-E_F)=552$~mV. Symmetric bias: $\mu_{L,R}=E_F \pm V/2$.}
\label{fig:device}
\end{figure*}

%% Fig 2: Methodology flowchart
\setcounter{figure}{1}
\begin{figure}[!t]
\centering
\includegraphics[width=1.1\columnwidth,height=17.5cm,keepaspectratio]{fig3_flowchart.png}
\caption{Rank-1 projected NEGF--SCBA methodology. Molecular coupling is rank-1 in transverse modes for any grid size, reducing 3D SCBA to one 1D projected equation. Anderson mixing (DIIS, $M=8$, $\beta=0.3$) ensures convergence with current conservation $<0.2\%$. $d^2I/dV^2$ computed by numerical double-differentiation of converged $I(V)$.}
\label{fig:flowchart}
\end{figure}

\begin{table}[!t]
\centering
\caption{Candidate RTD material stacks for room-temperature IETS. A viable sensor requires CBO $\sim$0.5~eV (NDR voltage in the 50--400~meV vibrational band), demonstrated NDR, and BEOL compatibility ($<400^\circ$C deposition). Only ZnO/Mg$_{0.3}$Zn$_{0.7}$O satisfies all criteria.}
\label{tab:stacks}
\renewcommand{\arraystretch}{1.05}
\footnotesize
\begin{tabular}{@{}lccccc@{}}
\toprule
Material stack & $\Delta E_c$ (eV) & $m^*/m_0$ & NDR & BEOL & Ref. \\
\midrule
AlGaAs/GaAs & 0.23 & 0.067 & Yes & No & \cite{sollner1983} \\
GaN/AlN & 1.8 & 0.20 & Yes & No & \cite{bayram2010} \\
In$_2$O$_3$/$\kappa$-Ga$_2$O$_3$ & 0.45 & 0.30 & No & Yes & \cite{kuang2021}\cite{zhu2024} \\
\textbf{ZnO/Mg$_{0.3}$Zn$_{0.7}$O} & \textbf{0.47} & \textbf{0.28} & \textbf{Yes} & \textbf{Yes} & \cite{krishnamoorthy2002}$^\ddagger$ \\
\bottomrule
\end{tabular}
\\\footnotesize $^\ddagger$ Krishnamoorthy \textit{et al.} used $x=0.2$; this simulation uses $x=0.3$ \\(CBO scaling $1.57x$ \cite{yin2017}).
\end{table}

%% Fig 3: (a) IV  (b) d2I/dV2  (c) DeltaD
\setcounter{figure}{2}
\begin{figure}[!t]
\centering
\begin{subfigure}[b]{0.48\columnwidth}
  \includegraphics[width=\textwidth,height=3.0cm]{fig4_IV.png}
  \caption{I--V at 300~K.}
  \label{fig:iv}
\end{subfigure}
\hfill
\begin{subfigure}[b]{0.48\columnwidth}
  \includegraphics[width=\textwidth,height=3.0cm]{fig5_d2IdV2.png}
  \caption{$d^2I/dV^2$ at 300~K.}
  \label{fig:d2i}
\end{subfigure}
\\[2pt]
\begin{subfigure}[b]{\columnwidth}
  \includegraphics[width=\textwidth,height=3.2cm]{fig7_deltaD.png}
  \caption{Discrimination metric $\Delta D(V)$ at 300~K.}
  \label{fig:deltaD}
\end{subfigure}
\caption{Simulation results at 300~K for Baseline, Mol\_A ($\hbar\omega=100$~meV), Mol\_B ($\hbar\omega=180$~meV), and Mol\_AB on a $10\,\mu$m~$\times$~$10\,\mu$m ZnO/Mg$_{0.3}$Zn$_{0.7}$O RTD (51 bias points, 16~mV spacing). \textbf{(a)}~I--V: NDR at $\sim$560~mV, Baseline peak $\sim$69~nA (PVR~$\approx$12). \textbf{(b)}~Numerical $d^2I/dV^2$: molecular modes produce species-specific deviations. \textbf{(c)}~Discrimination metric $\Delta D = d^2I_{\mathrm{mol}}/dV^2 - d^2I_{\mathrm{BL}}/dV^2$. Mol\_AB is approximately the superposition of Mol\_A and Mol\_B corrections, consistent with rank-1 additive structure.}
\label{fig:results}
\end{figure}

%% Fig 4: Temperature dependence
\setcounter{figure}{3}
\begin{figure}[!t]
\centering
\includegraphics[width=\columnwidth]{temp.png}
\caption{Temperature dependence (10--300~K) for Mol\_A ($\hbar\omega=100$~meV). \textbf{(a)}~I--V: peak current rises from $\sim$69~nA (300~K) to $\sim$110~nA (10~K) as Fermi smearing decreases; NDR sharpens. \textbf{(b)}~$d^2I/dV^2$; inset magnifies the 0--260~mV region. The molecular feature near 100~mV narrows at low $T$ and persists at 300~K, confirming that RTD energy filtering preserves IETS features beyond the thermal broadening limit $3k_BT$.}
\label{fig:temp}
\end{figure}

%% Per IEEE Sensors 2026 instructions, references must be confined to page 4.
\clearpage

\begin{thebibliography}{99}
%% Bibitems are ordered to match first-citation order in the body text,
%% so the printed numbers [1]..[N] run sequentially through the document.

\bibitem{hansma1977}
P.~K.~Hansma,
``Inelastic electron tunneling,''
\textit{Phys.\ Rep.}, vol.~30, no.~2, pp.~145--206, 1977.

\bibitem{hansma1982}
P.~K.~Hansma, Ed.,
\textit{Tunneling Spectroscopy: Capabilities, Applications, and New Techniques},
1st ed. New York, NY, USA: Springer, 1982.

\bibitem{hipps1993}
K.~W.~Hipps and U.~Mazur,
``Inelastic electron tunneling: An alternative molecular spectroscopy,''
\textit{J.\ Phys.\ Chem.}, vol.~97, no.~30, pp.~7803--7814, 1993.

\bibitem{reed2008}
M.~A.~Reed,
``Inelastic electron tunneling spectroscopy,''
\textit{Mater.\ Today}, vol.~11, no.~11, pp.~46--50, Nov.\ 2008.

\bibitem{jeong2015}
H.~Jeong \textit{et al.},
``Investigation of inelastic electron tunneling spectra of metal--molecule--metal junctions fabricated using direct metal transfer method,''
\textit{Appl.\ Phys.\ Lett.}, vol.~106, no.~6, art.\ no.~063110, Feb.\ 2015, doi: 10.1063/1.4908185.

\bibitem{liu2010}
Z.~Liu \textit{et al.},
``Inelastic electron tunneling spectroscopy study of ultrathin Al$_2$O$_3$--TiO$_2$ dielectric stack on Si,''
\textit{Appl.\ Phys.\ Lett.}, vol.~97, no.~20, art.\ no.~202905, Nov.\ 2010, doi: 10.1063/1.3518478.

\bibitem{kim2011}
E.~J.~Kim, M.~Shandalov, K.~C.~Saraswat, and P.~C.~McIntyre,
``Inelastic electron tunneling study of crystallization effects and defect energies in hafnium oxide gate dielectrics,''
\textit{Appl.\ Phys.\ Lett.}, vol.~98, no.~3, art.\ no.~032108, Jan.\ 2011, doi: 10.1063/1.3527977.

\bibitem{wang2004}
W.~Wang,
``Electrical characterization of self-assembled monolayers,''
Ph.D.\ dissertation, Yale Univ., New Haven, CT, USA, 2004.

\bibitem{lauhon2001}
L.~Lauhon and W.~Ho,
``Effects of temperature and other experimental variables on single molecule vibrational spectroscopy with the scanning tunneling microscope,''
\textit{Rev.\ Sci.\ Instrum.}, vol.~72, no.~1, pp.~216--223, Jan.\ 2001.

\bibitem{turin1996}
L.~Turin,
``A spectroscopic mechanism for primary olfactory reception,''
\textit{Chem.\ Senses}, vol.~21, no.~6, pp.~773--791, Dec.\ 1996, doi: 10.1093/chemse/21.6.773.

% \bibitem{brookes2007}
% J.~C.~Brookes, F.~Hartoutsiou, A.~P.~Horsfield, and A.~M.~Stoneham,
% ``Could humans recognize odor by phonon assisted tunneling?,''
% \textit{Phys.\ Rev.\ Lett.}, vol.~98, no.~3, art.\ no.~038101, Jan.\ 2007, doi: 10.1103/PhysRevLett.98.038101.

\bibitem{patil2018}
A.~Patil, D.~Saha, and S.~Ganguly,
``A quantum biomimetic electronic nose sensor,''
\textit{Sci.\ Rep.}, vol.~8, art.\ no.~128, 2018, doi: 10.1038/s41598-017-18346-2.



\bibitem{bommisetty2008}
V.~Bommisetty \textit{et al.},
``Gas sensing based on inelastic electron tunneling spectroscopy,''
\textit{IEEE Sens.\ J.}, vol.~8, no.~6, pp.~983--988, Jun.\ 2008.

\bibitem{patil2013mit}
P.~Patil,
``Design and fabrication of electron energy filters for room temperature inelastic electron tunneling spectroscopy,''
M.S.\ thesis, Massachusetts Inst.\ Technol., Cambridge, MA, USA, 2013.

\bibitem{li2024}
M.~Li, C.~S.~Cucinotta, and A.~P.~Horsfield,
``A computational model for a molecular chemical sensor,''
\textit{Nanoscale}, vol.~16, no.~10, pp.~5334--5342, 2024, doi: 10.1039/d3nr05900f.

\bibitem{pandey2021}
N.~Pandey, D.~Pal, D.~Saha, and S.~Ganguly,
``Vibration-based biomimetic odor classification,''
\textit{Sci.\ Rep.}, vol.~11, art.\ no.~11389, May 2021, doi: 10.1038/s41598-021-90592-x.

\bibitem{rock2008}
C. R\"ock, N. Barsan, and U. Weimar, ``Electronic nose: current status and future trends,'' \emph{Chem. Rev.}, vol.\ 108, no.\ 2, pp.\ 705--725, 2008.

\bibitem{jaklevic1966}
R.~C.~Jaklevic and J.~Lambe,
``Molecular vibration spectra by electron tunneling,''
\textit{Phys.\ Rev.\ Lett.}, vol.~17, no.~22, pp.~1139--1140, Nov.\ 1966.

\bibitem{ngabonziza2021}
P.~Ngabonziza, Y.~Wang, P.~A.~van~Aken, J.~Maier, and J.~Mannhart,
``Inelastic electron tunneling spectroscopy at high temperatures,''
\textit{Adv.\ Mater.}, vol.~33, no.~8, art.\ no.~2007299, Feb.\ 2021, doi: 10.1002/adma.202007299.

\bibitem{krishnamoorthy2002}
S.~Krishnamoorthy, A.~A.~Iliadis, A.~Inumpudi, S.~Choopun, R.~D.~Vispute, and T.~Venkatesan,
``Observation of resonant tunneling action in ZnO/Zn$_{0.8}$Mg$_{0.2}$O devices,''
\textit{Solid-State Electron.}, vol.~46, no.~10, pp.~1633--1637, 2002, doi: 10.1016/S0038-1101(02)00117-X.

\bibitem{kesarwani2022}
S.~Kesarwani, S.~Misra, D.~Saha, M.~L.~Della~Rocca, I.~Roy, S.~Ganguly, and A.~Mahajan,
``Simplified inelastic electron tunneling spectroscopy based on low-noise derivatives,''
\textit{Sci.\ Rep.}, vol.~12, art.\ no.~19216, Nov.\ 2022, doi: 10.1038/s41598-022-21302-4.

\bibitem{datta_textbook}
S.~Datta,
\textit{Quantum Transport: Atom to Transistor}.
Cambridge, U.K.: Cambridge Univ.\ Press, 2005.

\bibitem{lake1992}
R.~Lake and S.~Datta,
``Nonequilibrium Green's-function method applied to double-barrier resonant-tunneling diodes,''
\textit{Phys.\ Rev.\ B}, vol.~45, no.~12, pp.~6670--6685, Mar.\ 1992.

\bibitem{lake1997}
R.~Lake, G.~Klimeck, R.~C.~Bowen, and D.~Jovanovic,
``Single and multiband modeling of quantum electron transport through layered semiconductor devices,''
\textit{J.\ Appl.\ Phys.}, vol.~81, no.~12, pp.~7845--7869, Jun.\ 1997.

\bibitem{akkala2011}
A.~G.~Akkala,
``NEGF simulation of electron transport in resonant tunneling and resonant interband tunneling diodes,''
M.S.\ thesis, School of Electrical and Computer Engineering, Purdue Univ., West Lafayette, IN, USA, Dec.\ 2011.

\bibitem{anderson1965}
D.~G.~Anderson,
``Iterative procedures for nonlinear integral equations,''
\textit{J.\ ACM}, vol.~12, no.~4, pp.~547--560, Oct.\ 1965.

\bibitem{sollner1983}
T.~C.~L.~G.~Sollner, W.~D.~Goodhue, P.~E.~Tannenwald, C.~D.~Parker, and D.~D.~Peck,
``Resonant tunneling through quantum wells at frequencies up to 2.5~THz,''
\textit{Appl.\ Phys.\ Lett.}, vol.~43, no.~6, pp.~588--590, Sep.\ 1983, doi: 10.1063/1.94434.

\bibitem{bayram2010}
C.~Bayram, Z.~Vashaei, and M.~Razeghi,
``AlN/GaN double-barrier resonant tunneling diodes grown by metal-organic chemical vapor deposition,''
\textit{Appl.\ Phys.\ Lett.}, vol.~96, no.~4, art.\ no.~042103, Jan.\ 2010, doi: 10.1063/1.3294633.

\bibitem{kuang2021}
Y.~Kuang, X.~Chen, T.~Ma, Q.~Du, Y.~Zhang, J.~Hao, F.~Ren, B.~Liu, S.~Zhu, S.~Gu, R.~Zhang, Y.~Zheng, and J.~Ye,
``Band alignment and enhanced interfacial conductivity manipulated by polarization in a surfactant-mediated grown $\kappa$-Ga$_2$O$_3$/In$_2$O$_3$ heterostructure,''
\textit{ACS Appl.\ Electron.\ Mater.}, vol.~3, no.~2, pp.~795--803, 2021, doi: 10.1021/acsaelm.0c00947.

\bibitem{zhu2024}
J.~Zhu, J.~Li, S.~Ju, L.~Lu, S.~Zhang, and X.~Wang,
``Ultrathin In$_2$O$_3$ thin-film transistors deposited from trimethylindium and ozone,''
\textit{Nanotechnology}, vol.~35, no.~43, art.\ no.~435205, 2024, doi: 10.1088/1361-6528/ad6993.

\bibitem{yin2017}
H.~Yin, J.~Chen, Y.~Wang, J.~Wang, and H.~Guo,
``Composition dependent band offsets of ZnO and its ternary alloys,''
\textit{Sci.\ Rep.}, vol.~7, art.\ no.~41567, 2017, doi: 10.1038/srep41567.

\end{thebibliography}

\end{document}
